\section{Numerical experiments}\label{sec:experiments}

We conduct two numerical experiments to evaluate
Algorithm~\ref{alg:prox_admm}. The first concerns one-dimensional
piecewise-constant signal denoising, for which the first-order
difference matrix has full row rank.
The second concerns total variation denoising on the MNIST dataset,
for which the two-dimensional difference matrix does not have full row rank. 
All experiments were run in Python on a Windows~11 PC with a
13th-generation Intel Core i5 processor and 16~GB of RAM.
The data-generation code and competing-algorithm settings are given in
the Appendix.

\subsection{Piecewise-constant signal denoising}\label{Signal}

We recover a piecewise-constant signal \(b\in\mathbb{R}^n\)
from a noisy observation \(\hat b\in\mathbb{R}^n\).
Table~\ref{tab:models_algorithms} lists the tested models and corresponding
solvers. Let \(D\in\mathbb{R}^{(n-1)\times n}\) be the first-order difference
matrix with \(-1\) on the main diagonal and \(1\) on the
superdiagonal.

\begin{table}[h]
\centering
\footnotesize
\caption{Optimization models and corresponding solvers}
\label{tab:models_algorithms}
%\renewcommand{\arraystretch}{1}
\begin{tabular}{@{} l c c @{}}
\toprule
\textbf{Model} & \textbf{Problem formulation} & \textbf{Solver} \\
\midrule
\(\ell_1\)-\(C_{\ell_1}\)
&
$
\min_{x,y}\;
\|x-\hat b\|_1/n
+\lambda_p\|Dx-y\|
+\lambda_0\Phi(y)
$
&
SAP-ADMM\textsuperscript{H}/SAP-ADMM
\\
\midrule
\(\ell_1\)-\(\ell_{1/2}\)
&
$
\min_x\;
\|x-\hat b\|_1/n
+\lambda_{\ell_{1/2}}\|Dx\|_{1/2}^{1/2}
$
&
SDCAM\textsuperscript{1}~\cite{L2019A}
\\
\midrule
\(\ell_1\)-\(\ell_1\)
&
$
\min_{x,y}\;
\|x-\hat b\|_1/n
+\lambda_p\|Dx-y\|
+\lambda_{\ell_1}\|y\|_1
$
&
SAP-ADMM\textsuperscript{1}
\\
\midrule
\(\ell_2\)-\(\ell_{1/2}\)
&
$
\min_x\;
\|x-\hat b\|^2/2
+\lambda_{\ell_{1/2}}\|Dx\|_{1/2}^{1/2}
$
&
SDCAM\textsuperscript{2}~\cite{L2019A}
\\
\midrule
\(\ell_2\)-\(\ell_0\)
&
$
\min_x\;
\|x-\hat b\|^2/2
+\lambda_{\ell_0}\|Dx\|_0
$
&
pADMM~\cite{bot2020The}
\\
\bottomrule
\end{tabular}

\vspace{2pt}
\begin{minipage}{0.98\linewidth}
\footnotesize
\textbf{Note.}
The abbreviation \(\ell_1\)-\(C_{\ell_1}\) denotes the
\(\ell_1\)-capped-\(\ell_1\) model. Motivated by the acceleration observed when incorporating a restart
strategy into SAP-ADMM with $\alpha=2$ and $t=1$, we consider
SAP-ADMM\textsuperscript{H}, the Halpern variant with $\alpha=2$, $t=1$,
and a single restart at iteration $1000$. In comparison, SAP-ADMM
uses $\alpha=15$ and $t=1.5$. SDCAM\textsuperscript{1} and SDCAM\textsuperscript{2}
correspond to the $\ell_1$-$\ell_{1/2}$ and
$\ell_2$-$\ell_{1/2}$ models, respectively.
\end{minipage}
\end{table}


We set \(n=1000\). The segment lengths and signal levels are sampled from \([50,150]\) and \([-5,10]\), respectively, with adjacent levels differing by more than \(2\). The observation \(\hat b\) is generated by adding Gaussian noise with standard deviation \(0.5\) and impulse noise with probability \(\pi \in \{0,0.05,\ldots,0.20\}\). Each corrupted entry is perturbed uniformly on \([-5,5]\). We perform $50$ independent trials for each $\pi$.
Jump-location recovery is measured by the $F_1$ score, which is computed as
\[
F_1
=
\frac{
2|\mathcal{S}_{\mathrm{true}}\cap\mathcal{S}_{\mathrm{est}}|
}{
2|\mathcal{S}_{\mathrm{true}}\cap\mathcal{S}_{\mathrm{est}}|
+
|\mathcal{S}_{\mathrm{est}}\setminus\mathcal{S}_{\mathrm{true}}|
+
|\mathcal{S}_{\mathrm{true}}\setminus\mathcal{S}_{\mathrm{est}}|
}.
\]
Clearly, $F_1=1$ corresponds to exact support recovery.
Set
\(\varepsilon_{\mathrm{supp}}=10^{-6}\), and define the set of true jump locations \(\mathcal{S}_{\mathrm{true}}
=\bigl\{i:|(Db)_i|>\varepsilon_{\mathrm{supp}}\bigr\}\). For SAP-ADMM, SAP-ADMM\textsuperscript{H}, and
SAP-ADMM\textsuperscript{1}, the estimated jump locations are
\(\mathcal{S}_{\mathrm{est}}
=\{i:|\bar y_i|>\varepsilon_{\mathrm{supp}}\}\),
where \(\bar y\) is returned by the corresponding algorithm
at termination.
For the remaining algorithms, \(\mathcal{S}_{\mathrm{est}}\)
consists of the indices for which the absolute value of the
corresponding output component exceeds
\(\varepsilon_{\mathrm{supp}}\).

For the \(\ell_1\)-\(C_{\ell_1}\) model, we set
\(\lambda_0=0.016\) and \(\lambda_p=500L_f\), where
\(L_f=1/\sqrt{n}\), ensuring that Assumption~\ref{asp} holds.
For SAP-ADMM and SAP-ADMM\textsuperscript{H}, we set
\(N_{\mathrm{max}}=500\), \(\rho=2/n\), and
\(\beta=6\rho+10^{-8}\).
The initial points are chosen as \(x^0=\hat b\),
\(y^0=\mathbf{0}\), \(p^0=x^0\), \(q^0=Dx^0-y^0\),
\(\eta^0=\mathbf{0}\), and \(\mu^0=\mathbf{0}\).
To improve support identification, we set \(\nu_0=2\) and update
\[
\nu_k
=\max\left\{
\nu_{k-1}\gamma_\nu,\;
0.99\,\lambda_0/(3\lambda_p)
\right\},
\qquad k\geq1,
\]
where \(\gamma_\nu=0.999\) for \(k\leq1500\) and
\(\gamma_\nu=0.95\) otherwise. Here, \(k\) counts effective
iterations, excluding rejected acceleration trials.
After finitely many iterations, \(\nu_k\) reaches and remains
at \(\nu_{\min}=0.99\,\lambda_0/(3\lambda_p)\), which satisfies
Assumption~\ref{asn_1}.
Define
\[
\Delta_k
:=
\max\bigl\{
\|\bar p^k-p^k\|,
\|\bar q^k-q^k\|,
\|\bar x^k-x^k\|,
\|\bar y^k-y^k\|,
\|\bar\eta^k-\eta^k\|,
\|\bar\mu^k-\mu^k\|
\bigr\}.
\]
{\color{blue}
Note that $\Delta_k=0$ implies $R(\bar w^k)=\mathbf{0}$.}
We terminate SAP-ADMM and SAP-ADMM\textsuperscript{H} when
\(\Delta_k<2\times10^{-4}\), provided that \(\nu_k\) has reached
\(\nu_{\min}\) and at least \(50\) additional effective iterations
have been performed, or when the iteration count reaches
\(2\times10^4\).
The parameter settings and stopping criteria used for the
competing algorithms are provided in the Appendix.
 
\begin{figure}[htbp]
    \centering
    \includegraphics[width=\textwidth]{F1_Time_Efficiency.pdf}
    \caption{Performance comparison of six algorithms under
    different impulse noise probabilities:
    (a) average \(F_1\) score with standard deviation;
    (b) average computation time;
    (c) efficiency ratio.}
    \label{fig:f1_time_efficiency}
\end{figure}


Figure~\ref{fig:f1_time_efficiency} reports the average \(F_1\)
score and computation time, as well as the efficiency ratio, defined
as the average \(F_1\) score divided by the average computation time.
The results in Figure~\ref{fig:f1_time_efficiency}(a) show that
SAP-ADMM and SAP-ADMM\textsuperscript{H} applied to the
\(\ell_1\)-\(C_{\ell_1}\) model, and SDCAM\textsuperscript{1}
applied to the \(\ell_1\)-\(\ell_{1/2}\) model, consistently
attain average \(F_1\) scores above \(0.8\).
In contrast, SAP-ADMM\textsuperscript{1} applied to the
\(\ell_1\)-\(\ell_1\) model is less accurate in identifying jumps.
This comparison illustrates the advantage of nonconvex penalties
for sparse jump recovery.
SDCAM\textsuperscript{2} and pADMM, which solve the
\(\ell_2\)-\(\ell_{1/2}\) and \(\ell_2\)-\(\ell_0\) models,
respectively, perform well without impulse noise but deteriorate
rapidly as \(\pi\) increases, compared with the algorithms
solving models with the \(\ell_1\) loss.
The computation times in Figure~\ref{fig:f1_time_efficiency}(b) indicate
that pADMM, SAP-ADMM, SAP-ADMM\textsuperscript{H}, and
SAP-ADMM\textsuperscript{1} incur relatively low computational
costs when solving their respective models. \textcolor{blue}{Figure~\ref{fig:f1_time_efficiency}(c) shows that SAP-ADMM
outperforms SAP-ADMM\textsuperscript{H} and all the other
compared algorithms when applied to their respective models in terms of computational efficiency,
indicating the advantage of the safeguarded acceleration strategy.}
Specifically, SAP-ADMM achieves a higher efficiency ratio than
SAP-ADMM\textsuperscript{1}, which in turn outperforms
SAP-ADMM\textsuperscript{H}.

\begin{table}[htbp]
\centering
\footnotesize
\caption{Average MSE and computation time for piecewise-constant signal denoising}
\label{tab:mse_time_comparison}
\renewcommand{\arraystretch}{1.10}
\setlength{\tabcolsep}{3.5pt}
\begin{tabular}{@{}lccccc@{}}
\toprule
\multirow{2}{*}{Algorithm}
& \multicolumn{5}{c}{Impulse noise probability \(\pi\)} \\
\cmidrule(lr){2-6}
& \(0.00\) & \(0.05\) & \(0.10\) & \(0.15\) & \(0.20\) \\
\midrule
SAP-ADMM
& \(0.0059/\mathbf{0.3453}\)
& \(\mathbf{0.0114}/\mathbf{0.3583}\)
& \(0.0189/0.3472\)
& \(\mathbf{0.0308}/0.3427\)
& \(0.0374/0.3690\) \\
SAP-ADMM\textsuperscript{H}
& \(0.0057/0.9521\)
& \(0.0115/0.9165\)
& \(\mathbf{0.0187}/0.9530\)
& \(0.0309/0.9643\)
& \(\mathbf{0.0369}/0.9392\) \\
SDCAM\textsuperscript{1}
& \(0.0157/6.2683\)
& \(0.0220/6.6608\)
& \(0.0287/6.4973\)
& \(0.0397/6.7519\)
& \(0.0449/6.9264\) \\
SAP-ADMM\textsuperscript{1}
& \(0.2256/0.3757\)
& \(0.2769/0.3596\)
& \(0.2230/\mathbf{0.3262}\)
& \(0.2075/\mathbf{0.3400}\)
& \(0.2019/\mathbf{0.3340}\) \\
SDCAM\textsuperscript{2}
& \(\mathbf{0.0033}/3.4328\)
& \(0.0174/3.6816\)
& \(0.0422/4.1424\)
& \(0.0810/4.3617\)
& \(0.1394/4.3245\) \\
pADMM
& \(0.0036/0.5546\)
& \(0.3933/0.5472\)
& \(0.7680/0.4289\)
& \(1.1339/0.3852\)
& \(1.5655/0.3681\) \\
\bottomrule
\end{tabular}

\vspace{2pt}
\begin{minipage}{0.96\linewidth}
\footnotesize
\textbf{Note.}
Each entry is reported as MSE/computation time (s).
\end{minipage}
\end{table}

Table~\ref{tab:mse_time_comparison} reports the average mean
squared error (MSE) and computation time.
The MSE is computed as \(\mathrm{MSE}=\|\bar x-b\|^2/n\),
where \(\bar x\) is the restored signal returned by the
corresponding algorithm at termination.
When \(\pi=0\), the \(\ell_2\)-\(\ell_{1/2}\) model
solved by SDCAM\textsuperscript{2} yields the lowest MSE.
For \(\pi\geq0.05\), the lowest MSE values are
obtained from the \(\ell_1\)-\(C_{\ell_1}\) model using
SAP-ADMM or SAP-ADMM\textsuperscript{H}.
Moreover, the \(\ell_1\)-\(C_{\ell_1}\) model solved by
SAP-ADMM and SAP-ADMM\textsuperscript{H}, together with the
\(\ell_1\)-\(\ell_{1/2}\) model solved by
SDCAM\textsuperscript{1}, consistently yields relatively low
MSE values across noise levels, illustrating the effectiveness
of the \(\ell_1\) loss combined with nonconvex penalties.
For the \(\ell_1\)-\(C_{\ell_1}\) model, SAP-ADMM substantially
reduces the computation time relative to SAP-ADMM\textsuperscript{H}.
Hence, we focus on SAP-ADMM
in all subsequent experiments.

\begin{figure}[htbp]
    \centering
    \includegraphics[width=\textwidth]{F1_Distribution.pdf}
    \caption{Distribution of \(F_1\) scores under different
    impulse noise probabilities.}
    \label{fig:f1_distribution}
\end{figure}

Figure~\ref{fig:f1_distribution} presents the distributions
of \(F_1\) scores over \(50\) trials.
SAP-ADMM\textsuperscript{1} applied to the
\(\ell_1\)-\(\ell_1\) model is omitted because it rarely
recovers the exact support.
SAP-ADMM and SDCAM\textsuperscript{1}, applied to the
\(\ell_1\)-\(C_{\ell_1}\) and
\(\ell_1\)-\(\ell_{1/2}\) models, respectively, attain
exact recovery in more than \(70\%\) of the trials for
\(\pi\leq0.05\) and in more than \(30\%\) of the trials
for \(\pi\in\{0.15,0.20\}\).
SDCAM\textsuperscript{2} and pADMM, applied to the
\(\ell_2\)-\(\ell_{1/2}\) and \(\ell_2\)-\(\ell_0\)
models, respectively, attain reliable exact recovery when
\(\pi=0\) but do not recover the exact support in any trial when
\(\pi\in\{0.15,0.20\}\).
For a direct comparison, Figure~\ref{fig:recovery} shows the
reconstructions obtained in one trial for \(\pi=0\) and \(0.20\),
using \(N_{\mathrm{max}}=2000\) for SAP-ADMM.
When \(\pi=0\), all compared algorithms except
SAP-ADMM\textsuperscript{1} closely recover the clean signal
when solving their respective models.
When \(\pi=0.20\), the \(\ell_1\)-\(C_{\ell_1}\) model
solved by SAP-ADMM and the \(\ell_1\)-\(\ell_{1/2}\) model
solved by SDCAM\textsuperscript{1} yield the most accurate
reconstructions, whereas those obtained from the
\(\ell_2\)-loss models using SDCAM\textsuperscript{2}
and pADMM retain noticeable noise.

\begin{figure}[htbp]
    \centering
    \includegraphics[width=\textwidth]{Figure2_Recovery_Comparison_FiveModels_acc_capl1_python.pdf}
    \caption{Signal recovery by the compared algorithms applied
    to their respective models under impulse noise probabilities
    of \(0\%\) and \(20\%\).}
    \label{fig:recovery}
\end{figure}

{\color{blue}
We further evaluate SAP-ADMM on the problem
\(
\min_{x,y}\;
\|Ax-\hat b\|_1/m
+\lambda_p\|Dx-y\|
+\lambda_0\Phi(y),
\)
where \(m=800\) and \(n=1000\).
In each trial, \(A\in\mathbb{R}^{m\times n}\) is generated
with independent standard Gaussian entries and normalized
so that \(\|A\|=1\).
The signal \(b\) and matrix \(D\) are constructed as before.
The observation \(\hat b\) is generated from \(Ab\), with Gaussian standard deviation
\(0.2\) and impulse noise probability \(0.05\).
We set \(\lambda_0=0.0018\), \(\lambda_p=500/\sqrt{m}\),
and \(\rho=0.6/n\), and initialize
\(x^0=A^\top\hat b\) and \(p^0=Ax^0\). Here, \(\nu\) is updated as above, except that
\(\gamma_\nu=0.9995\) for \(k\leq4000\) and
\(\gamma_\nu=0.95\) thereafter.
The remaining parameters, initial points, and stopping
criterion for SAP-ADMM are unchanged.
We perform \(20\) independent trials.
Within each trial, the same problem instance is used for all
\(N_{\mathrm{max}}\in\{200,300,400,500,600\}\).

\begin{table}[htbp]
\color{blue}
\centering
\small
\caption{Average SAP-ADMM results for different values of \(N_{\mathrm{max}}\)}
\label{tab:nmax_general_A}
\renewcommand{\arraystretch}{1.05}
\setlength{\tabcolsep}{10pt}
\begin{tabular}{cccc}
\toprule
\(N_{\mathrm{max}}\)
& Average \(F_1\)
& Average iterations
& Average MSE \\
\midrule
200 & 1.0000 & 8239.95 & 0.00344 \\
300 & 1.0000 & 8239.95 & 0.00344 \\
400 & 1.0000 & 8239.95 & 0.00344 \\
500 & 1.0000 & 8239.95 & 0.00344 \\
600 & 1.0000 & 8239.95 & 0.00344 \\
\bottomrule
\end{tabular}
\end{table}

Table~\ref{tab:nmax_general_A} reports the average \(F_1\)
score, iteration count, and MSE over the \(20\) trials.
The average \(F_1\) score and MSE demonstrate the effectiveness
of SAP-ADMM in signal recovery when \(A\) is not an identity matrix.
The reported averages remain unchanged across all tested
values of \(N_{\mathrm{max}}\), indicating that SAP-ADMM
is relatively insensitive to this parameter under the
considered settings.
}

All these results demonstrate that SAP-ADMM achieves a favorable balance
between reconstruction accuracy and computational efficiency
for piecewise-constant signal denoising.

\subsection{MNIST total variation denoising}

We consider total variation denoising of \(28\times28\)
MNIST images~\cite{LeCun1998}.
Let \(b\in\mathbb{R}^n\), with \(n=n_1n_2\), represent the
columnwise vectorization of a clean \(n_1\times n_2\) image,
and let \(\hat b\) be its noisy observation.
Here, \(D=[D_h^\top,D_v^\top]^\top\), with
\(D_h=L_{n_2}\otimes I_{n_1}\) and
\(D_v=I_{n_2}\otimes L_{n_1}\), where \(\otimes\) denotes the
Kronecker product and \(L_{n_i}\in\mathbb{R}^{n_i\times n_i}\)
is the first-order difference matrix with entries
\((L_{n_i})_{j,j}=-1\) and \((L_{n_i})_{j,j+1}=1\)
for \(j=1,\ldots,n_i-1\), and all other entries equal to zero.
By \(D\in\mathbb{R}^{2n_1n_2\times n_1n_2}\) and
\(\operatorname{rank}(D)\leq n_1n_2<2n_1n_2\), it follows
that \(D\) does not have full row rank.
Since the convergence of pADMM requires \(D\) to
have full row rank, we test all models in Table~\ref{tab:models_algorithms} except the $\ell_2$-$\ell_0$ model solved by pADMM.

For each digit \(0,\ldots,9\), we select one image and generate
\(20\) noisy observations by replacing \(10\%\) of its pixels
with independent samples from the uniform distribution on
\([0,1]\), adding Gaussian noise with standard deviation
\(0.2\), and clipping the resulting values to \([0,1]\).
For the \(\ell_1\)-\(C_{\ell_1}\) model, we set
\(\lambda_0=8\times10^{-5}\) and \(\lambda_p=2L_f\).
For SAP-ADMM, we set \(\rho=0.5/n\) and
\(\beta=10\rho+10^{-8}\), with \(n=28\times28\).
Starting from \(\nu_0=0.1\), we update \(\nu_k\) according to
\[
\nu_k
=
\max\left\{
\nu_0\,0.99^k,\;
0.999\,\lambda_0/(3\lambda_p)
\right\}.
\]
Here, \(\nu_{\min}=0.999\,\lambda_0/(3\lambda_p)\).
The stopping criterion is \(\Delta_k<4\times10^{-4}\), subject
to the same requirement of at least \(50\) additional effective
iterations after \(\nu_k\) reaches \(\nu_{\min}\).
The remaining parameters, initial points, and maximum iteration
number are the same as those used in \cref{Signal}.
The settings of the competing algorithms are listed in
the Appendix.

We evaluate restoration quality using the peak signal-to-noise
ratio (PSNR) for pixelwise reconstruction accuracy, the structural
similarity index measure (SSIM)~\cite{Wang2004} for structural
similarity, and the gradient magnitude similarity deviation
(GMSD)~\cite{Xue2014} for spatial variation in local gradient
magnitude similarity.
Higher PSNR and SSIM values and lower GMSD values indicate
better restoration quality.
The calculation formulas are provided in the Appendix.

Table~\ref{tab:mnist_four_models_compact} reports the average
PSNR, computation time, SSIM, and GMSD.
The \(\ell_1\)-\(C_{\ell_1}\) model solved by SAP-ADMM
yields the best average PSNR, SSIM, and GMSD for every digit.
Compared with the \(\ell_1\)-\(\ell_1\) model solved by
SAP-ADMM\textsuperscript{1}, it improves average PSNR by
approximately \(3.22\%\), increases SSIM by \(14.45\%\), and
reduces GMSD by \(27.43\%\).
When solving their respective models, SAP-ADMM requires,
on average, more computation time than SAP-ADMM\textsuperscript{1} and
SDCAM\textsuperscript{2}, but substantially less than
SDCAM\textsuperscript{1}.

\begin{table}[htbp]
\centering
\caption{Comparison of four algorithms applied to their respective
models for MNIST image denoising}
\label{tab:mnist_four_models_compact}
\scriptsize
\setlength{\tabcolsep}{3.6pt}
\renewcommand{\arraystretch}{1.08}
\begin{tabular}{@{}c l c c c c@{}}
\toprule
Digit & Metric
& SAP-ADMM
& SDCAM\textsuperscript{1}
& SAP-ADMM\textsuperscript{1}
& SDCAM\textsuperscript{2} \\
\midrule
\multirow{2}{*}{0}
& \(P\uparrow/T\downarrow\) & \(\mathbf{18.05}/0.1890\) & \(17.34/0.7503\) & \(17.77/\mathbf{0.0583}\) & \(15.01/0.2787\) \\
& \(S\uparrow/G\downarrow\) & \(\mathbf{0.629}/\mathbf{0.1958}\) & \(0.535/0.2176\) & \(0.561/0.2556\) & \(0.392/0.2459\) \\
\midrule
\multirow{2}{*}{1}
& \(P\uparrow/T\downarrow\) & \(\mathbf{21.25}/0.3414\) & \(19.86/0.7839\) & \(20.48/\mathbf{0.0623}\) & \(16.38/0.2066\) \\
& \(S\uparrow/G\downarrow\) & \(\mathbf{0.484}/\mathbf{0.1491}\) & \(0.365/0.2064\) & \(0.388/0.2783\) & \(0.222/0.2169\) \\
\midrule
\multirow{2}{*}{2}
& \(P\uparrow/T\downarrow\) & \(\mathbf{17.57}/0.1768\) & \(16.83/0.8604\) & \(17.25/\mathbf{0.0573}\) & \(14.67/0.2682\) \\
& \(S\uparrow/G\downarrow\) & \(\mathbf{0.609}/\mathbf{0.2286}\) & \(0.525/0.2378\) & \(0.559/0.2622\) & \(0.396/0.2752\) \\
\midrule
\multirow{2}{*}{3}
& \(P\uparrow/T\downarrow\) & \(\mathbf{18.92}/0.4346\) & \(17.92/0.8890\) & \(18.52/\mathbf{0.0599}\) & \(15.28/0.2203\) \\
& \(S\uparrow/G\downarrow\) & \(\mathbf{0.557}/\mathbf{0.1941}\) & \(0.475/0.2190\) & \(0.497/0.2513\) & \(0.328/0.2566\) \\
\midrule
\multirow{2}{*}{4}
& \(P\uparrow/T\downarrow\) & \(\mathbf{18.92}/0.2863\) & \(18.55/0.8587\) & \(18.12/\mathbf{0.0575}\) & \(15.60/0.2411\) \\
& \(S\uparrow/G\downarrow\) & \(\mathbf{0.435}/\mathbf{0.1981}\) & \(0.381/0.2128\) & \(0.391/0.2729\) & \(0.253/0.2410\) \\
\midrule
\multirow{2}{*}{5}
& \(P\uparrow/T\downarrow\) & \(\mathbf{18.31}/0.4181\) & \(17.44/0.9173\) & \(17.57/\mathbf{0.0570}\) & \(15.19/0.2886\) \\
& \(S\uparrow/G\downarrow\) & \(\mathbf{0.550}/\mathbf{0.2063}\) & \(0.463/0.2218\) & \(0.473/0.2602\) & \(0.326/0.2445\) \\
\midrule
\multirow{2}{*}{6}
& \(P\uparrow/T\downarrow\) & \(\mathbf{19.44}/0.1536\) & \(18.28/0.7825\) & \(18.79/\mathbf{0.0589}\) & \(15.47/0.1928\) \\
& \(S\uparrow/G\downarrow\) & \(\mathbf{0.521}/\mathbf{0.1761}\) & \(0.426/0.2111\) & \(0.458/0.2705\) & \(0.304/0.2312\) \\
\midrule
\multirow{2}{*}{7}
& \(P\uparrow/T\downarrow\) & \(\mathbf{19.19}/0.2871\) & \(18.06/0.8059\) & \(18.35/\mathbf{0.0561}\) & \(15.55/0.2751\) \\
& \(S\uparrow/G\downarrow\) & \(\mathbf{0.531}/\mathbf{0.1824}\) & \(0.422/0.2099\) & \(0.440/0.2657\) & \(0.289/0.2268\) \\
\midrule
\multirow{2}{*}{8}
& \(P\uparrow/T\downarrow\) & \(\mathbf{18.15}/0.3975\) & \(17.21/0.7685\) & \(17.41/\mathbf{0.0557}\) & \(14.65/0.2733\) \\
& \(S\uparrow/G\downarrow\) & \(\mathbf{0.531}/\mathbf{0.1994}\) & \(0.451/0.2258\) & \(0.470/0.2660\) & \(0.323/0.2602\) \\
\midrule
\multirow{2}{*}{9}
& \(P\uparrow/T\downarrow\) & \(\mathbf{18.93}/0.5348\) & \(18.00/0.8203\) & \(18.57/\mathbf{0.0565}\) & \(15.37/0.2510\) \\
& \(S\uparrow/G\downarrow\) & \(\mathbf{0.533}/\mathbf{0.1897}\) & \(0.436/0.2114\) & \(0.463/0.2624\) & \(0.299/0.2402\) \\
\midrule
\multirow{2}{*}{Avg.}
& \(P\uparrow/T\downarrow\) & \(\mathbf{18.87}/0.3219\) & \(17.95/0.8237\) & \(18.28/\mathbf{0.0579}\) & \(15.32/0.2496\) \\
& \(S\uparrow/G\downarrow\) & \(\mathbf{0.538}/\mathbf{0.1920}\) & \(0.448/0.2174\) & \(0.470/0.2645\) & \(0.313/0.2438\) \\
\bottomrule
\end{tabular}

\vspace{2pt}
\begin{minipage}{\linewidth}
\footnotesize
\textbf{Note.}
\(P\), \(T\), \(S\), and \(G\) denote PSNR, computation time (s),
SSIM, and GMSD, respectively.
The symbols \(\uparrow\) and \(\downarrow\) indicate that
larger and smaller values are preferable, respectively.
\end{minipage}
\end{table}

\begin{figure}[h]
    \centering
    \includegraphics[width=0.7\textwidth]{mnist_four_model_restoration_python.pdf}
    \caption{Original images, noisy observations, and restorations
    obtained by the four compared algorithms applied to their
    respective models.}
    \label{fig:mnist}
\end{figure}

For a visual comparison, Figure~\ref{fig:mnist} presents the
restored images from one trial.
The \(\ell_1\)-\(C_{\ell_1}\) model solved by SAP-ADMM yields
cleaner backgrounds and sharper digit contours than the
other tested models solved by their corresponding algorithms.
The \(\ell_1\)-\(\ell_1\) model solved by
SAP-ADMM\textsuperscript{1} produces the most blurred digits,
consistent with its largest average GMSD value in
Table~\ref{tab:mnist_four_models_compact}.
These results indicate that, even when \(D\) does not have full
row rank, SAP-ADMM remains numerically effective in solving the
\(\ell_1\)-\(C_{\ell_1}\) model and performs well in image denoising.
